\documentclass[11pt,a4paper]{article}

\usepackage[T1]{fontenc}
\usepackage[utf8]{inputenc}
\usepackage{lmodern}
\usepackage[margin=1in]{geometry}
\usepackage{amsmath,amssymb,amsfonts,amsthm,bm}
\usepackage{graphicx}
\usepackage{booktabs}
\usepackage{longtable}
\usepackage{tabularx}
\usepackage{array}
\usepackage{enumitem}
\usepackage{caption}
\usepackage{subcaption}
\usepackage{xcolor}
\usepackage{fancyhdr}
\usepackage{microtype}
\usepackage[round,authoryear]{natbib}
\usepackage[hidelinks]{hyperref}
\usepackage{verbatim}
\usepackage{tikz}
\usetikzlibrary{arrows.meta,positioning,fit,backgrounds}

\newcommand{\paperversion}{1.0.0}
\newcommand{\codecommit}{\texttt{ef13675}}
\newcommand{\evidencetag}{\texttt{final-\allowbreak audit-\allowbreak 2026-10-\allowbreak 31-\allowbreak dpa}}
\newcommand{\F}{\mathbb{F}}
\newcommand{\bits}{\{0,1\}}
\newcommand{\Rec}{\mathsf{Rec}}
\newcommand{\Share}{\mathsf{Share}}
\newcommand{\HW}{\mathsf{HW}}
\newcommand{\Sbox}{\mathsf{SBOX}}
\newcommand{\code}[1]{\texttt{#1}}
\newcommand{\Adv}{\mathcal{A}}
\newcommand{\E}{\mathbb{E}}
\newcommand{\Prb}{\Pr}
\newcommand{\negl}{\mathsf{negl}}

\theoremstyle{plain}
\newtheorem{theorem}{Theorem}[section]
\newtheorem{lemma}[theorem]{Lemma}
\newtheorem{proposition}[theorem]{Proposition}
\newtheorem{corollary}[theorem]{Corollary}
\theoremstyle{definition}
\newtheorem{definition}[theorem]{Definition}
\newtheorem{assumption}[theorem]{Assumption}
\newtheorem{claim}[theorem]{Claim}
\theoremstyle{remark}
\newtheorem{remark}[theorem]{Remark}

\newcounter{algblock}
\newenvironment{algorithmblock}[1]{%
  \par\medskip\refstepcounter{algblock}%
  \noindent\textbf{Algorithm \thealgblock} --- #1\par\nopagebreak
  \noindent\begin{minipage}{\linewidth}\small
}{%
  \end{minipage}\par\medskip
}

\pagestyle{fancy}
\fancyhf{}
\fancyhead[L]{\small SiCaS Yellow Paper}
\fancyhead[R]{\small Version \paperversion}
\fancyfoot[C]{\small\thepage}
\renewcommand{\headrulewidth}{0.4pt}

\hypersetup{
  pdftitle={SiCaS Yellow Paper: Formal Specification and Verification Obligations},
  pdfauthor={SiCaS Contributors},
  pdfsubject={Formal specification of side-channel attenuation and leakage assessment},
  pdfkeywords={formal specification, Boolean masking, TVLA, CPA, DPA, ChaCha20, Rust, verification},
  colorlinks=true,
  linkcolor=blue!45!black,
  citecolor=blue!45!black,
  urlcolor=blue!45!black
}

\title{%
  \vspace{-1.2cm}
  {\LARGE\bfseries SiCaS}\\[0.35cm]
  {\Large Yellow Paper}\\[0.25cm]
  {\large Formal Specification and Verification Obligations}
}
\author{SiCaS Contributors}
\date{Version \paperversion \quad 31 October 2026}

\setlength{\emergencystretch}{2em}

\begin{document}

\maketitle
\thispagestyle{empty}

\begin{center}
\begin{tabular}{@{}ll@{}}
\toprule
Document & Yellow paper \\
Version & \paperversion \\
Date & 31 October 2026 \\
Code commit & \codecommit \\
Release build commit & \texttt{123eb69} \\
Evidence tag & \evidencetag \\
Public label & SAL-2 locally \\
Production status & Not certified, not audited, not hardware-validated \\
License & GPL-3.0-only \\
\bottomrule
\end{tabular}
\end{center}

\begin{abstract}
This yellow paper specifies the functional behaviour, data formats, security
definitions, and verification obligations of SiCaS, a safe-Rust workspace for
first-order side-channel attenuation and leakage assessment. It defines the two-share
Boolean masking model, the masked operations, the masked ChaCha20 reference
integration, the 64-block bitsliced batch path, the trace format, the TVLA/CPA/DPA
algorithms, the timing and leakage-model checks, and the release manifest. It states
the correctness theorems and the first-order security argument, lists the assumptions
under which those statements hold, and maps each public claim to a proof obligation and
an evidence artifact. The specification is deliberately explicit about what is not
proven: physical leakage, higher-order attacks, fault injection, template attacks,
production zeroization, and independent audit remain outside the current claim set. The
document is a specification and verification plan, not a substitute for independent
review.
\end{abstract}

\vspace{0.3cm}
\noindent\textbf{Keywords:} formal specification, Boolean masking, TVLA, CPA, DPA,
ChaCha20, Rust, proof obligations, assurance.

\newpage
\tableofcontents
\newpage

\section{Scope and relationship to the white paper}

This document is the formal companion to the SiCaS white paper. The white paper
explains the motivation, architecture, and evaluation at a technical-management level.
The yellow paper specifies:

\begin{enumerate}[leftmargin=1.5em]
  \item the data types and invariants of the core masking library;
  \item the functional behaviour of the masked operations;
  \item the ChaCha20 reference integration and the bitsliced batch path;
  \item the trace format and the TVLA, CPA, and DPA algorithms;
  \item the security definitions and the assumptions behind the first-order claim;
  \item the verification obligations and the evidence mapping;
  \item the release, provenance, and reproducibility formats.
\end{enumerate}

The specification is written against code commit \codecommit{} and evidence tag
\evidencetag. It is a \emph{specification of the current implementation}, not a
specification of an ideal cryptographic library. Where the implementation has a known
limitation, the limitation is stated.

\section{Notation and conventions}

\begin{description}[leftmargin=1.5em]
  \item[\(\bits\)] The set \(\{0,1\}\).
  \item[\(\bits^n\)] The set of \(n\)-bit strings.
  \item[\(\F_2\)] The field with two elements.
  \item[\(\F_{2^8}\)] The AES field, represented as bytes with the AES reduction
        polynomial.
  \item[\(\oplus\)] Bitwise XOR.
  \item[\(\wedge\)] Bitwise AND.
  \item[\(\neg\)] Bitwise complement.
  \item[\(\ll\)] Left shift.
  \item[\(\gg\)] Right shift.
  \item[\(\lll\)] Rotate left.
  \item[\(\ggg\)] Rotate right.
  \item[\(\HW(x)\)] Hamming weight of \(x\).
  \item[\(x_b\)] Bit \(b\) of \(x\), with \(x_0\) the least significant bit.
  \item[{\([n]\)}] The set \(\{0,\ldots,n-1\}\).
  \item[\(\gets\)] Uniform random sampling.
  \item[\(\gets_{\$}\)] Uniform random sampling from a finite set.
  \item[{\(\Prb[E]\)}] Probability of event \(E\).
  \item[{\(\E[X]\)}] Expectation of random variable \(X\).
  \item[\(\negl(\lambda)\)] A negligible function in security parameter \(\lambda\).
\end{description}

All arithmetic on \(n\)-bit words is modulo \(2^n\) unless stated otherwise. Byte order
in the trace format is little-endian. The notation \(u16_{\mathrm{le}}\) and
\(u32_{\mathrm{le}}\) denotes little-endian encoding.

\section{Formal model}\label{sec:model}

\subsection{Values, shares, and reconstruction}

\begin{definition}[Two-share Boolean sharing]\label{def:sharing}
For \(x\in\bits^n\), a two-share Boolean sharing is a pair
\((s_0,s_1)\in\bits^n\times\bits^n\) such that
\[
  \Rec(s_0,s_1) := s_0\oplus s_1 = x.
\]
The sharing is \emph{fresh} if it is generated as
\[
  r\gets_{\$}\bits^n,\qquad s_1:=r,\qquad s_0:=x\oplus r.
\]
\end{definition}

\begin{definition}[Masked operation]\label{def:masked-op}
Let \(f:(\bits^n)^m\to\bits^n\) be an \(m\)-ary operation. A masked implementation
\(f^\ast\) takes \(m\) fresh sharings and returns a sharing
\((c_0,c_1)\) such that
\[
  \Rec(c_0,c_1) = f(x_1,\ldots,x_m)
\]
for all inputs \(x_1,\ldots,x_m\), with overwhelming probability over the fresh masks.
\end{definition}

\begin{definition}[Declassification boundary]\label{def:declass}
A declassification boundary is an operation that intentionally reveals a value that
was protected by a secret type or a sharing. Examples are \code{Secret::expose\_secret},
\code{Masked::unmask}, and the ChaCha20 keystream output.
\end{definition}

\subsection{Leakage functions and first-order game}

\begin{definition}[Leakage function]\label{def:leakage}
A leakage function is a map \(L:\bits^m\to\mathcal{Y}\) for some output set
\(\mathcal{Y}\). A first-order single-share leakage function depends on at most one
share of a sharing.
\end{definition}

\begin{definition}[First-order single-share game]\label{def:game}
Let \(\lambda\) be a security parameter and \(n=n(\lambda)\). The game
\(\mathsf{Exp}^{\mathsf{fo}}_{\Adv,\mathcal{L}}(n)\) proceeds as follows:
\begin{enumerate}[leftmargin=1.5em]
  \item The challenger samples \(x\gets_{\$}\bits^n\), \(r\gets_{\$}\bits^n\), and
        sets \(s_0:=x\oplus r\), \(s_1:=r\).
  \item The adversary \(\Adv\) chooses a leakage function
        \(L\in\mathcal{L}\) and a bit \(b\gets_{\$}\bits\).
  \item The challenger returns \(L(s_b)\) to \(\Adv\).
  \item \(\Adv\) outputs a guess \(b'\).
\end{enumerate}
The advantage is
\[
  \mathsf{Adv}^{\mathsf{fo}}_{\Adv,\mathcal{L}}(n)
  := \left|\Prb[b'=b]-\frac{1}{2}\right|.
\]
\end{definition}

\begin{theorem}[Ideal single-share security]\label{thm:ideal}
For any \(x\in\bits^n\), any uniform independent mask \(r\), and any leakage function
\(L\) that depends on only one share, the distribution of \(L(s_b)\) is independent of
\(x\). Consequently, for all adversaries \(\Adv\),
\[
  \mathsf{Adv}^{\mathsf{fo}}_{\Adv,\mathcal{L}}(n)=0
\]
in the ideal model.
\end{theorem}

\begin{proof}
Each share is uniform on \(\bits^n\) and independent of \(x\): \(s_1=r\) is uniform,
and \(s_0=x\oplus r\) is uniform for every fixed \(x\). Therefore the marginal
distribution of either share is the uniform distribution, which does not depend on
\(x\). Any function of a single share has a distribution that depends only on that
share's distribution.
\end{proof}

\begin{remark}[What the theorem does not say]
Theorem~\ref{thm:ideal} models an ideal single-share leakage function. It does not
model glitches, transitions, coupling between shares, register reuse, or a physical
probe that observes a function of both shares. Those effects are outside the current
claim and are the reason hardware validation is required.
\end{remark}

\subsection{Assumptions}

\begin{assumption}[Fresh masks]\label{asm:fresh}
Every mask is sampled uniformly from \(\bits^n\), is independent of the secret and of
all other masks, and is used at most once.
\end{assumption}

\begin{assumption}[Entropy]\label{asm:entropy}
The entropy source is a secure CSPRNG when it reports success. A failure is reported as
an error and is never replaced by a weak fallback.
\end{assumption}

\begin{assumption}[Constant-time compilation]\label{asm:ct}
The compiler preserves the data-independent control flow and memory-access pattern of
the constant-time gadgets on the evaluated target.
\end{assumption}

\begin{assumption}[Leakage model]\label{asm:model}
The assessment uses a leakage model that is appropriate for the target. A passing test
under an inappropriate model is not evidence of absence of leakage.
\end{assumption}

\begin{assumption}[No faults or invasive access]\label{asm:fault}
The adversary does not inject faults, modify the device, or perform invasive probing.
\end{assumption}

\section{Core data types and invariants}\label{sec:types}

\subsection{\code{Choice}}

\code{Choice} is a one-bit value used for branchless selection. It exposes
\code{from\_bool}, \code{as\_u8}, and mask/selection helpers. The invariant is that a
\code{Choice} is represented by a value in \(\{0,1\}\) and that no operation on it
branches on a secret.

\subsection{\code{Secret<T>}}

\code{Secret<T>} wraps a value and deliberately does not implement \code{PartialEq},
\code{Ord}, \code{Hash}, \code{Index}, \code{Clone}, or \code{Copy}. The only ways to
read the value are \code{expose\_secret} and \code{into\_inner}, both of which are
declassification boundaries. The type does not guarantee zeroization.

\subsection{\code{Masked<T>}}

A \code{Masked<T>} value is a pair of shares \((s_0,s_1)\). The invariant is
\[
  s_0\oplus s_1 = x
\]
where \(x\) is the logical value. The shares are not exposed except through
\code{shares}, \code{unmask}, and the explicit analysis boundary. A fresh sharing is
constructed by \code{Masked::mask(value, mask)} with a caller-supplied uniform mask, or
by \code{Masked::from\_parts(s\_0,s\_1)} when the caller has already constructed a valid
sharing.

\subsection{\code{EntropySource}}

\begin{definition}[Entropy source contract]\label{def:entropy}
An \code{EntropySource} implementation must either fill the entire destination buffer
with cryptographically secure random bytes and return \code{Ok(())}, or return
\code{Err(EntropyError::Unavailable)} without claiming success. Partial fill followed
by success is a contract violation.
\end{definition}

\subsection{Invariants}

\begin{longtable}{@{}p{0.08\linewidth}p{0.50\linewidth}p{0.36\linewidth}@{}}
\caption{Core invariants.}\label{tab:invariants}\\
\toprule
ID & Invariant & Enforcement \\
\midrule
\endfirsthead
\toprule
ID & Invariant & Enforcement \\
\midrule
\endhead
I1 & Reconstruction: \(s_0\oplus s_1=x\) for every \code{Masked<T>}. & Tests, Kani \\
I2 & Freshness: every mask is uniform, independent, and used once. & API contract, review \\
I3 & Share independence: each share is uniform under a fresh mask. & Kani, statistical tests \\
I4 & Constant time: no secret-dependent branch or index in the core gadgets. & Assembly, dudect \\
I5 & Fail closed: entropy errors propagate as \code{Err}. & Tests, code review \\
I6 & No unsafe code in core, reference, platform, or harness. & \code{forbid(unsafe\_code)} \\
I7 & Parser bounds: trace dimensions are checked before allocation. & Tests, fuzzing \\
I8 & Release reproducibility: identical artifact hash across independent builds. & Two-build compare \\
\bottomrule
\end{longtable}

\section{Constant-time primitives}\label{sec:ct}

\subsection{Comparison}

For \(n\in\{8,16,32,64,128\}\), \code{ct::eq\_n(a,b)} returns a \code{Choice} that is
\(1\) iff \(a=b\). The implementation is branchless:
\[
  \code{eq}_n(a,b) = \neg\bigvee_{i=0}^{n-1}(a_i\oplus b_i).
\]
The \code{is\_zero\_n} family is the same predicate with \(b=0\).

\subsection{Selection and swap}

For a public \code{Choice} \(c\),
\[
  \code{select}_n(c,a,b) = (a\wedge \code{mask}(c))\oplus(b\wedge\neg\code{mask}(c)),
\]
where \(\code{mask}(c)\) is the all-ones word when \(c=1\) and zero otherwise. The swap
operation uses the same mask to exchange two words without a branch.

\subsection{Table lookup}

\code{lookup\_u8(table, index)} scans all 256 entries and selects the entry whose index
matches, without a secret-dependent memory access. The table must have exactly 256
entries. The implementation is constant-time with respect to the index.

\begin{theorem}[Constant-time lookup correctness]\label{thm:lookup}
For every table \(T\in(\bits^8)^{256}\) and every index \(i\in\bits^8\),
\(\code{lookup\_u8}(T,i)=T[i]\).
\end{theorem}

\begin{proof}
The scan accumulates \(T[j]\) exactly when \(j=i\), using a branchless equality mask.
All other terms are masked out. The result is \(T[i]\).
\end{proof}

\section{Masked arithmetic}\label{sec:masked}

\subsection{XOR, complement, shifts, and rotations}

These operations are share-wise:
\begin{align}
  (a_0,a_1)\oplus(b_0,b_1) &= (a_0\oplus b_0,\;a_1\oplus b_1),\\
  \neg(a_0,a_1) &= (\neg a_0,\;a_1),\\
  (a_0,a_1)\lll k &= (a_0\lll k,\;a_1\lll k),\\
  (a_0,a_1)\ggg k &= (a_0\ggg k,\;a_1\ggg k),\\
  (a_0,a_1)\ll k &= (a_0\ll k,\;a_1\ll k),\\
  (a_0,a_1)\gg k &= (a_0\gg k,\;a_1\gg k).
\end{align}
For shifts, amounts greater than or equal to the word width produce zero. The
implementation uses checked shifts and returns zero for out-of-range amounts.

\subsection{Refresh}

\[
  \code{refresh}((a_0,a_1),r) = (a_0\oplus r,\;a_1\oplus r).
\]
The logical value is unchanged, and the new shares are randomized by \(r\).

\subsection{AND}

\begin{algorithmblock}{Masked AND}\label{alg:and}
\begin{verbatim}
Input:  shares (a0, a1), (b0, b1), fresh uniform mask r
Output: shares (c0, c1)

c0 <- (a0 & b0) ^ r
c1 <- (a0 & b1) ^ (a1 & b0) ^ (a1 & b1) ^ r
return (c0, c1)
\end{verbatim}
\end{algorithmblock}

\begin{theorem}[Masked AND correctness]\label{thm:and}
For all \(a,b\in\bits^n\) and all valid sharings, the output of
Algorithm~\ref{alg:and} satisfies \(c_0\oplus c_1=a\wedge b\).
\end{theorem}

\begin{proof}
Expand \((a_0\oplus a_1)\wedge(b_0\oplus b_1)\) into the four products. The two
occurrences of \(r\) cancel under XOR.
\end{proof}

\begin{theorem}[Masked AND first-order freshness]\label{thm:and-fresh}
If \(r\) is uniform and independent of the inputs, then each output share is uniform
and independent of the logical value.
\end{theorem}

\begin{proof}
\(c_0\) is the XOR of a fixed value with the uniform mask \(r\), hence uniform. \(c_1\)
is also the XOR of a fixed value with \(r\), hence uniform. The pair reconstructs
\(a\wedge b\) by Theorem~\ref{thm:and}.
\end{proof}

\subsection{Masked addition}

\begin{algorithmblock}{Masked 32-bit addition}\label{alg:add}
\begin{verbatim}
Input:  shares (a0, a1), (b0, b1), entropy source E
Output: shares (c0, c1)

result <- (0, 0)
carry  <- (0, 0)
for i = 0 .. 31:
    a <- bit(a, i)
    b <- bit(b, i)
    sum <- a XOR b XOR carry
    result <- result XOR sum
    if i == 31: break
    first  <- AND(a, b, E.mask_u32())
    second <- AND(a, carry, E.mask_u32())
    third  <- AND(b, carry, E.mask_u32())
    carry  <- (first XOR second XOR third) << 1
return result
\end{verbatim}
\end{algorithmblock}

\begin{theorem}[Masked addition correctness]\label{thm:add}
For all \(a,b\in\bits^{32}\), the output of Algorithm~\ref{alg:add} satisfies
\(\Rec(c_0,c_1)=(a+b)\bmod 2^{32}\).
\end{theorem}

\begin{proof}[Proof sketch]
The reconstructed sum bit at position \(i\) is \(a_i\oplus b_i\oplus c_i\), and the
reconstructed carry is
\((a_i\wedge b_i)\oplus(a_i\wedge c_i)\oplus(b_i\wedge c_i)\), shifted left by one.
This is the standard full-adder recurrence. Induction on \(i\) from \(c_0=0\) gives
the 32-bit sum modulo \(2^{32}\). Each masked AND preserves the reconstructed value by
Theorem~\ref{thm:and}, so the masked implementation computes the same function.
\end{proof}

\begin{remark}[Fresh masks]
The correctness proof does not require fresh masks; the security argument does. Reusing
a mask across ANDs invalidates Assumption~\ref{asm:fresh} and the first-order claim.
\end{remark}

\subsection{Selection, equality, and bit extraction}

Selection selects each share independently with a public \code{Choice}. Equality
compares \(a_0\oplus b_0\) with \(a_1\oplus b_1\) in constant time. Bit extraction masks
both shares with \(1\ll i\), returning zero for \(i\ge n\). These operations preserve
the reconstruction invariant.

\section{ChaCha20 specification}\label{sec:chacha}

\subsection{State and constants}

The ChaCha20 state is sixteen 32-bit words:
\[
  \text{state} = (c_0,c_1,c_2,c_3,\;k_0,\ldots,k_7,\;counter,\;n_0,n_1,n_2),
\]
where the constants are
\[
  (c_0,c_1,c_2,c_3) = (\texttt{0x61707865},\texttt{0x3320646e},
  \texttt{0x79622d32},\texttt{0x6b206574}),
\]
the key is 256 bits, the counter is 32 bits, and the nonce is 96 bits.

\subsection{Quarter round}

\begin{algorithmblock}{ChaCha20 quarter round}
\begin{verbatim}
Input:  state words a, b, c, d
Output: updated state words

a <- a + b
d <- (d XOR a) <<< 16
c <- c + d
b <- (b XOR c) <<< 12
a <- a + b
d <- (d XOR a) <<< 8
c <- c + d
b <- (b XOR c) <<< 7
\end{verbatim}
\end{algorithmblock}

\subsection{Block function}

\begin{algorithmblock}{ChaCha20 block}
\begin{verbatim}
Input:  key (32 bytes), nonce (12 bytes), counter (u32)
Output: 64-byte keystream block

state <- constants || key || counter || nonce
initial <- state
for 10 double rounds:
    quarter_round(state, 0, 4, 8, 12)
    quarter_round(state, 1, 5, 9, 13)
    quarter_round(state, 2, 6, 10, 14)
    quarter_round(state, 3, 7, 11, 15)
    quarter_round(state, 0, 5, 10, 15)
    quarter_round(state, 1, 6, 11, 12)
    quarter_round(state, 2, 7, 8, 13)
    quarter_round(state, 3, 4, 9, 14)
for i = 0 .. 15:
    state[i] <- state[i] + initial[i]
return little_endian_bytes(state)
\end{verbatim}
\end{algorithmblock}

\subsection{Masked scalar implementation}

The scalar implementation stores the state as sixteen \code{Masked<u32>} words. The
initial state is masked with fresh randomness. For each block, the working state and
the initial state used by the final addition are independently refreshed. The quarter
round uses masked addition, XOR, and rotation. The final addition uses the masked
adder. The keystream output is unmasked only at the declassification boundary.

\subsection{Bitsliced batch implementation}

The batch implementation processes \(B=64\) blocks. For each of the 16 state words and
each of the 32 bit positions, a \code{Masked<u64>} plane holds one bit from each of the
64 lanes. The quarter round operates on planes. The masked adder is a ripple-carry
adder over the 32 bit planes:
\[
  \text{carry}_{i+1} =
  (a_i\wedge b_i)\oplus(a_i\wedge \text{carry}_i)\oplus(b_i\wedge \text{carry}_i),
\]
with a fresh mask for each AND. After the final addition, the planes are transposed
into lane-major output.

\begin{algorithmblock}{Bitsliced masked addition}\label{alg:add-planes}
\begin{verbatim}
Input:  planes a[0..31], b[0..31], entropy source E
Output: planes result[0..31]

carry <- zero plane
for i = 0 .. 31:
    xor <- a[i] XOR b[i]
    result[i] <- xor XOR carry
    if i == 31: break
    generate   <- AND(a[i], b[i], E.mask_u64())
    propagated <- AND(xor, carry, E.mask_u64())
    carry      <- generate XOR propagated
return result
\end{verbatim}
\end{algorithmblock}

\begin{remark}
Algorithm~\ref{alg:add-planes} uses the generate/propagate form of the full adder.
It is equivalent to the three-AND form in Algorithm~\ref{alg:add} after XOR
simplification, but it uses two masked ANDs per bit instead of three.
\end{remark}

\subsection{Transpose}

\begin{algorithmblock}{64x64 bit-matrix transpose}
\begin{verbatim}
Input:  32 planes, each a u64 row
Output: 64 u32 lane words

rows[0..31] <- unmask(planes[0..31])
for width in [32, 16, 8, 4, 2, 1]:
    mask <- stage_mask(width)
    for row in 0 .. 63:
        if row & width == 0:
            swap <- (rows[row] XOR (rows[row+width] << width)) & mask
            rows[row] <- rows[row] XOR swap
            rows[row+width] <- rows[row+width] XOR (swap >> width)
return rows[0..63] truncated to u32
\end{verbatim}
\end{algorithmblock}

The transpose is a standard butterfly network. It is applied to the unmasked output
planes; the unmasking is a declassification boundary at the keystream output. The
transpose is validated by differential tests against the scalar and unmasked
references.

\section{Harness specification}\label{sec:harness}

\subsection{Trace format}

\begin{verbatim}
TraceFile ::= Magic Version SamplesPerTrace TraceCount Sample*
Magic            ::= "SiCaSTR1"                    (8 bytes)
Version          ::= u16le(1)
SamplesPerTrace  ::= u32le
TraceCount       ::= u32le
Sample           ::= i16le
\end{verbatim}

The parser enforces:
\begin{enumerate}[leftmargin=1.5em]
  \item magic equals \code{SiCaSTR1};
  \item version equals 1;
  \item \(1\le \text{trace\_count}\le 10^6\);
  \item \(1\le \text{samples\_per\_trace}\le 10^6\);
  \item \(\text{trace\_count}\cdot\text{samples\_per\_trace}\le 10^8\);
  \item the file contains exactly the declared number of samples.
\end{enumerate}

\subsection{Preprocessing}

\begin{description}[leftmargin=1.5em]
  \item[Trim.] \(T'_{i,j}=T_{i,j+s}\) for \(j\in[0,e-s)\).
  \item[Filter.] \(T'_{i,j}=\lfloor\frac{1}{|W_j|}\sum_{k\in W_j}T_{i,k}\rfloor\),
        where \(W_j=[\max(0,j-r),\min(n-1,j+r)]\) and \(r=\lfloor w/2\rfloor\).
  \item[Align.] Choose \(d^\ast=\arg\max_{d\in[-m,m]}\rho(R,T_i,d)\), then shift
        \(T_i\) by \(d^\ast\) and zero-fill.
\end{description}

\subsection{TVLA}

\begin{algorithmblock}{Fixed-vs-random TVLA}
\begin{verbatim}
Input:  fixed traces F[0..n0-1], random traces R[0..n1-1], threshold t
Output: report

for sample j = 0 .. samples-1:
    compute means, variances, and Welch t_j
    track max |t_j|, its sample, and the count above t
return report
\end{verbatim}
\end{algorithmblock}

The statistic is
\[
  t_j = \frac{\bar{F}_j-\bar{R}_j}
             {\sqrt{s_F^2/n_0+s_R^2/n_1}},
\]
with \(s^2\) the sample variance. The implementation clamps negative variances to zero
and returns zero when the denominator is zero.

\subsection{CPA}

\begin{algorithmblock}{CPA}
\begin{verbatim}
Input:  traces T[0..n-1][0..s-1], plaintexts p[0..n-1], model M
Output: report

for key k = 0 .. 255:
    model[i] <- M(p[i], k)
    for sample j = 0 .. s-1:
        rho[k][j] <- Pearson(model, T[:,j])
    best[k] <- max_j |rho[k][j]|
return argmax_k best[k]
\end{verbatim}
\end{algorithmblock}

\begin{lemma}[Hamming-weight complement symmetry]\label{lem:hw}
For all \(x,k\in\bits^8\),
\[
  \HW(x\oplus k) + \HW(x\oplus k\oplus \texttt{0xFF}) = 8.
\]
Therefore the Hamming-weight model has perfectly negatively correlated predictions for
\(k\) and \(k\oplus\texttt{0xFF}\), and the tie-break returns the lower key.
\end{lemma}

\begin{proof}
Complementing all eight bits of a byte maps a byte of weight \(w\) to a byte of weight
\(8-w\). The sum is 8.
\end{proof}

\subsection{DPA}

\begin{algorithmblock}{AES S-box DPA}
\begin{verbatim}
Input:  traces T[0..n-1][0..s-1], plaintexts p[0..n-1]
Output: report

for key k = 0 .. 255:
    for bit b = 0 .. 7:
        for sample j = 0 .. s-1:
            S0 <- {i : SBOX[p[i] XOR k]_b = 0}
            S1 <- {i : SBOX[p[i] XOR k]_b = 1}
            if S0 or S1 is empty: continue
            delta <- mean(T[S1][j]) - mean(T[S0][j])
            track max |delta|, signed delta, key, bit, sample
return report
\end{verbatim}
\end{algorithmblock}

\begin{lemma}[Why the S-box matters]\label{lem:sbox}
For a single-bit leakage function \(L(p,k)=(p\oplus k)_b\), the absolute difference of
means is independent of the key bits other than \(b\). Consequently, a single-bit
partition of \(p\oplus k\) cannot distinguish the full key in the ideal model.
\end{lemma}

\begin{proof}[Proof sketch]
For fixed \(b\), \((p\oplus k)_b=p_b\oplus k_b\). Changing \(k_b\) swaps the two
partitions; changing any other key bit leaves the partition unchanged. Swapping the
partitions negates the difference but preserves its absolute value. Therefore all keys
with the same bit \(b\) have the same absolute difference, and keys with the other bit
value have the same absolute difference as well. The absolute value is independent of
the key.
\end{proof}

The AES S-box provides a nonlinear map from \(p\oplus k\) to the partition bit, so the
partition depends on all key bits. This is why the DPA implementation uses
\(\Sbox[p\oplus k]_b\) rather than \((p\oplus k)_b\).

\subsection{Timing and leakage model}

The timing command compares fixed and random classes with Welch's \(t\)-test. The
leakage-model command computes the Pearson correlation between each share and the
secret for masked byte, masked AND, and masked addition. A negative control must be
detected. The thresholds and sample counts are parameters and must be published with
the result.

\section{Security definitions and claims}\label{sec:claims}

\subsection{First-order security definition}

\begin{definition}[First-order security]\label{def:fo-secure}
An implementation is first-order secure for a leakage class \(\mathcal{L}\) if, for
every adversary \(\Adv\) and every \(L\in\mathcal{L}\), the advantage in
Definition~\ref{def:game} is negligible in the security parameter.
\end{definition}

The ideal argument of Theorem~\ref{thm:ideal} gives first-order security for
single-share leakage under Assumptions~\ref{asm:fresh} and~\ref{asm:entropy}. The
physical implementation is not proven first-order secure; it is designed to satisfy the
model and evaluated with local evidence.

\subsection{Claims}

\begin{longtable}{@{}p{0.06\linewidth}p{0.44\linewidth}p{0.42\linewidth}@{}}
\caption{Public claims and status.}\label{tab:claims}\\
\toprule
ID & Claim & Status \\
\midrule
\endfirsthead
\toprule
ID & Claim & Status \\
\midrule
\endhead
C1 & The safe API prevents accidental secret-dependent branching/indexing. &
Compile-fail tests pass. \\
C2 & Constant-time gadgets are data-independent on supported targets. &
Dudect and assembly evidence on x86\_64; other targets pending. \\
C3 & First-order Boolean masking reduces first-order DPA below the TVLA threshold for
the reference algorithm. &
Synthetic/model evidence; hardware TVLA pending. \\
C4 & Entropy failure is fail-closed. &
Local tests pass. \\
C5 & Zero runtime dependencies. &
Verified locally. \\
C6 & No unsafe code in core logic crates. &
\code{forbid(unsafe\_code)}. \\
C7 & Builds are reproducible. &
Two-build compare passes. \\
C8 & Critical invariants are model-checked. &
30 Kani harnesses pass. \\
C9 & v1.0 meets SAL-4. &
Pending independent verification. \\
C10 & Release artifacts are signed and traceable. &
Local signing dry run; production pending. \\
\bottomrule
\end{longtable}

\subsection{Proof obligations}

\begin{longtable}{@{}p{0.08\linewidth}p{0.46\linewidth}p{0.38\linewidth}@{}}
\caption{Proof obligations and evidence.}\label{tab:obligations}\\
\toprule
ID & Obligation & Evidence / status \\
\midrule
\endfirsthead
\toprule
ID & Obligation & Evidence / status \\
\midrule
\endhead
PO1 & \code{Choice} and \code{ct} gadgets are branchless. & Assembly matrix, dudect; pass on x86\_64. \\
PO2 & \code{Secret} prevents accidental equality/indexing/cloning. & Compile-fail doctests; pass. \\
PO3 & \code{Masked} reconstruction is correct. & Unit tests, Kani; pass. \\
PO4 & Masked AND is correct and fresh. & Kani, tests; pass. \\
PO5 & Masked addition is correct. & Kani, tests; pass. \\
PO6 & Refresh preserves the value and randomizes shares. & Kani, tests; pass. \\
PO7 & Entropy failure propagates. & Tests; pass. \\
PO8 & ChaCha20 matches RFC~8439. & Test vectors; pass. \\
PO9 & Bitsliced batch matches scalar/unmasked. & Differential tests; pass. \\
PO10 & Trace parser rejects malformed input. & Fuzzing, Miri; pass. \\
PO11 & TVLA/CPA/DPA are correct on synthetic data. & Unit tests; pass. \\
PO12 & No unsafe code in core/reference/platform/harness. & Lints, source audit; pass. \\
PO13 & Release artifact is reproducible. & Two-build compare; pass. \\
PO14 & Zeroization drop path erases bytes. & Tests, Miri, assembly; pass for the adapter. \\
PO15 & Performance hard ceiling is met. & Benchmark; 7.15$\times$; pass. \\
PO16 & External audit, red team, hardware, production signing. & Pending. \\
\bottomrule
\end{longtable}

\section{Verification methodology}\label{sec:verification}

\subsection{Testing}

The workspace uses unit tests, property-style randomized tests, differential tests,
compile-fail doctests, and integration tests. The core tests cover the constant-time
gadgets, secret handling, masking, refresh, reconstruction, and addition. The reference
tests cover RFC~8439 vectors, differential checks, lane independence, key-plane balance,
batch/scalar interoperability, counter wrap, and entropy failure. The harness tests
cover parsing, preprocessing, TVLA, CPA, DPA, timing, synthetic generation, and leakage
model checks.

\subsection{Kani}

Kani \citep{kani} is used for bounded model checking of the core gadgets. The 30
harnesses cover constant-time comparison, selection, swap, lookup, masking, refresh,
share reconstruction, addition, and share uniformity. The proofs are bounded to the
declared integer widths and the assumptions in the harness.

\subsection{Miri and sanitizers}

Miri \citep{miri} checks the Rust mid-level IR for undefined behaviour on the covered
paths. AddressSanitizer \citep{serebryany2012asan} checks native memory errors.
Valgrind \citep{nethercote2007valgrind} checks native memory errors and definite leaks
on five test binaries. The statistical, OS, and subprocess paths that Miri cannot
execute are documented and run natively.

\subsection{Fuzzing and mutation}

Five libFuzzer campaigns \citep{libfuzzer} cover masked operations, constant-time
lookup, trace parsing, ChaCha20, and the bitsliced batch path. Mutation testing
\citep{jia2011mutation} measures the strength of the test suite: the current workspace
run tests 630 mutants, catches 574, misses 0, and classifies 56 as unviable.

\subsection{Assembly and timing}

The assembly checker inspects the generated code for secret-dependent branches and
memory accesses across optimization levels, codegen units, panic strategies, and LTO
variants. The dudect-style timing tests \citep{reparaz2017dudect} compare fixed and
random classes. The results are local measurements, not a proof for every target.

\subsection{Release verification}

The release manifest validator checks the artifact hash, checksum file, provenance,
SBOM component graph, workspace versions, timestamp, and (optionally) the current
commit. The reproducible-build comparison builds the artifact twice in independent
target directories and compares the bytes. The local signing dry run signs and verifies
a detached signature with an ephemeral test key.

\section{Data formats}\label{sec:formats}

\subsection{Trace and plaintext files}

The trace format is defined in Section~\ref{sec:harness}. The plaintext file is a raw
byte sequence with one byte per trace. The trace count and plaintext count must match
for CPA and DPA.

\subsection{Reports}

The harness prints human-readable reports. The TVLA report contains \code{max\_abs\_t},
\code{max\_sample}, \code{exceeding\_samples}, \code{sample\_count},
\code{fixed\_traces}, \code{random\_traces}, and \code{threshold}. The CPA report
contains \code{best\_key}, \code{best\_correlation}, \code{best\_sample}, and a
per-key correlation vector. The DPA report contains \code{best\_key}, \code{best\_bit},
\code{best\_sample}, \code{best\_difference}, and \code{best\_signed\_difference}.

\subsection{Release manifest}

The release manifest directory contains:

\begin{description}[leftmargin=1.5em]
  \item[\code{sicas-harness}] The release binary.
  \item[\code{SHA256SUMS}] The SHA-256 checksum of the binary.
  \item[\code{provenance.json}] Project, commit, rustc, host, build command, artifact,
        SHA-256, and timestamp.
  \item[\code{sbom.cdx.json}] CycloneDX SBOM with workspace components and dependency
        edges.
  \item[\code{README.md}] Human-readable summary with commit, host, Rust version,
        artifact, and SHA-256.
\end{description}

\subsection{Evidence bundle}

The evidence bundle contains the proof summaries, mutation and coverage evidence,
cross-architecture results, release metadata, signing dry-run record, and a source
snapshot. The bundle has an \code{INDEX.md} and a \code{MANIFEST.sha256}; the manifest is
verified with \code{sha256sum -c MANIFEST.sha256}.

\section{Error model and fail-closed behaviour}\label{sec:errors}

\begin{description}[leftmargin=1.5em]
  \item[\code{EntropyError::Unavailable}] The entropy source failed or is not
        supported. The caller must not fall back to a weak source.
  \item[\code{TraceError}] Invalid magic, unsupported version, invalid dimensions,
        oversized input, length mismatch, or I/O error. The parser returns a typed
        error and does not panic.
  \item[\code{TvlaError}] Empty group, too few traces, dimension mismatch, or invalid
        threshold.
  \item[\code{CpaError}] Empty input, length mismatch, or too few traces.
  \item[\code{DpaError}] Empty input, length mismatch, or too few traces.
  \item[\code{PreprocessError}] Empty trace, invalid range, invalid window, invalid
        reference, invalid shift limit, or trace rebuild failure.
\end{description}

The fail-closed principle applies to entropy and to release validation. A missing or
invalid input produces an error, not a silent fallback.

\section{Performance model}\label{sec:performance}

\subsection{Complexity}

\begin{itemize}[leftmargin=1.5em]
  \item Scalar ChaCha20: \(O(1)\) per 64-byte block; the state is fixed-size.
  \item Bitsliced batch: \(O(B)\) work for \(B=64\) blocks, with \(O(1)\) state per
        lane and a \(64\times64\) transpose.
  \item TVLA: \(O(N\cdot S)\) time and \(O(S)\) auxiliary space.
  \item CPA: \(O(256\cdot N\cdot S)\) time and \(O(S)\) auxiliary space per key.
  \item DPA: \(O(256\cdot 8\cdot N\cdot S)\) time and \(O(S)\) auxiliary space.
\end{itemize}

\subsection{Measured results}

The final full proof measures a scalar masked overhead of approximately \(320\times\)
and a bitsliced batch overhead of \(7.15\times\) versus the unmasked reference. The
10$\times$ hard ceiling is met. The 5$\times$ target is open; reaching it requires SIMD
or a faster audited mask stream.

\section{Conformance and test vectors}\label{sec:conformance}

\subsection{RFC 8439}

The reference tests include the RFC~8439 ChaCha20 block test vector and the
encryption test vector. The masked scalar and bitsliced paths are checked against the
unmasked reference for randomized keys, nonces, and counters. The tests also cover
counter wrap and batch/scalar interoperability.

\subsection{Differential statistics}

The TVLA, CPA, and DPA statistics are validated against an independent Python model
over 200 deterministic cases. The leakage-model check includes a negative control that
must be detected.

\subsection{Release conformance}

The release manifest validator checks that the artifact hash, checksum file,
provenance, SBOM, workspace versions, timestamp, and commit agree. The reproducible
build comparison checks that two independent builds produce identical bytes.

\section{Versioning and change log}\label{sec:versioning}

The document and the code are versioned separately. The document version follows
semantic versioning: major changes to the specification, minor additions, and patch
clarifications. The code version is the workspace version \code{0.1.0} until the
external gates pass. Every proof run records the code commit and the evidence tag.

\small
\begin{longtable}{@{}p{0.12\linewidth}p{0.18\linewidth}p{0.24\linewidth}p{0.38\linewidth}@{}}
\caption{Yellow paper revision history.}\label{tab:history}\\
\toprule
Version & Date & Code / evidence & Changes \\
\midrule
\endfirsthead
\toprule
Version & Date & Code / evidence & Changes \\
\midrule
\endhead
0.1.0 & 2026-10-30 & initial tree & Core types, masking model, trace format. \\
0.9.0 & 2026-10-31 & \code{8c41e4b} & Full proof obligations, evaluation mapping, release format. \\
1.0.0 & 2026-10-31 & \code{ef13675}, \evidencetag & AES S-box DPA, GPL-3.0-only, final audit package. \\
\bottomrule
\end{longtable}

\appendix

\section{API index}

\begin{longtable}{@{}p{0.34\linewidth}p{0.58\linewidth}@{}}
\caption{Public API index (selected).}\label{tab:api}\\
\toprule
Item & Role \\
\midrule
\endfirsthead
\toprule
Item & Role \\
\midrule
\endhead
\code{Choice::from\_bool}, \code{as\_u8} & Branchless one-bit value. \\
\code{ct::eq\_*}, \code{select\_*}, \code{swap\_u32}, \code{lookup\_u8} &
Constant-time comparison, selection, swap, and lookup. \\
\code{Secret::new}, \code{expose\_secret}, \code{into\_inner} &
Opaque secret wrapper and declassification. \\
\code{Masked::mask}, \code{from\_parts}, \code{shares}, \code{unmask} &
Sharing construction and reconstruction. \\
\code{Masked::refresh}, \code{xor}, \code{and}, \code{add\_masked} &
Masked operations. \\
\code{EntropySource::fill\_bytes}, \code{mask\_u32}, \code{mask\_u64} &
Fail-closed entropy contract. \\
\code{MaskedChaCha20::new}, \code{keystream\_block}, \code{keystream\_blocks} &
Masked ChaCha20 scalar and batch. \\
\code{TraceSet::from\_traces}, \code{read\_from}, \code{write\_to} &
Trace storage and I/O. \\
\code{tvla}, \code{cpa\_hamming\_weight}, \code{cpa\_with\_model}, \code{dpa\_bit} &
Leakage analysis. \\
\code{preprocess::trim}, \code{moving\_average}, \code{align} &
Trace preprocessing. \\
\code{timing::dudect}, \code{leakage::check\_all} &
Timing and leakage-model checks. \\
\code{OsEntropy::new} & Safe OS entropy adapter. \\
\code{zeroize\_bytes}, \code{ZeroizingBytes} & Experimental zeroization boundary. \\
\bottomrule
\end{longtable}

\section{Conformance checklist}

\begin{enumerate}[leftmargin=1.5em]
  \item Build the workspace with the pinned toolchain.
  \item Run all workspace tests and doctests.
  \item Run the quick local proof.
  \item Run the full local proof for the release candidate.
  \item Verify the evidence bundle manifest.
  \item Verify the release manifest and artifact hash.
  \item Compare the release SHA-256 with the published value.
  \item Check the claim-to-evidence map for the claimed assurance level.
  \item Record every external gate that remains open.
\end{enumerate}

\bibliographystyle{plainnat}
\bibliography{../references}

\end{document}
